\chapter{Los carbonilos: adiciones nucleófilas, acetales y protección}\label{ch:b1:acetals-protection}

Disuélvase glucosa en agua y casi nada de ella permanece como el aldehído
de cadena abierta que dibujan los libros: la molécula se curva, su propio
grupo hidroxilo se adiciona a su aldehído y se cierra un anillo de seis
miembros. La misma adición de un alcohol a un grupo carbonilo, repetida
una vez más con un \omterm{def:b1:elementary-steps:catalyst}{catalizador} ácido, da los \omterm{def:b1:acetals-protection:hemiacetal}{acetales}, compuestos estables
frente a las bases y a los \omterm{def:b1:electronic-effects:nucleophile}{nucleófilos} más reactivos. Esa estabilidad es
una herramienta: un aldehído que un \omterm{def:b1:organomagnesium:organometallic}{reactivo de Grignard} destruiría puede
ocultarse en forma de \omterm{def:b1:acetals-protection:hemiacetal}{acetal}, atravesar la reacción y recuperarse al
final. Este capítulo estudia el grupo carbonilo como \omterm{def:b1:electronic-effects:nucleophile}{electrófilo}, las
adiciones de agua y de alcoholes, y la estrategia de la \omterm{def:b1:acetals-protection:protecting-group}{protección}.

\begin{recall}
Las adiciones de Grignard a los compuestos carbonílicos
(\cref{ch:b1:organomagnesium}); el \omterm{def:b1:extent-q-and-k:quotient}{cociente de reacción} $Q$ y la
evolución de un sistema mientras $Q \neq K$ (\cref{ch:b1:extent-q-and-k});
las \omterm{def:b1:electronic-effects:curly-arrow}{flechas curvas}, los \omterm{def:b1:electronic-effects:nucleophile}{nucleófilos} y los \omterm{def:b1:electronic-effects:nucleophile}{electrófilos}
(\cref{ch:b1:electronic-effects}). El volumen escolar (año 12) introdujo
la idea de proteger un grupo durante una síntesis.
\end{recall}

\section{El grupo carbonilo como electrófilo}

El enlace C=O está polarizado hacia el oxígeno (\cref{ch:b1:periodicity}),
y una \omterm{def:b1:lewis-resonance-vsepr:resonance}{forma resonante} \ce{C+-O-} sitúa una carga positiva completa en el
carbono: el carbono es \omterm{def:b1:electronic-effects:nucleophile}{electrófilo}, y los \omterm{def:b1:electronic-effects:nucleophile}{nucleófilos} lo atacan
perpendicularmente al plano del grupo. Los \omterm{def:b1:electronic-effects:nucleophile}{nucleófilos} fuertes (\omterm{def:b1:organomagnesium:organometallic}{reactivos
de Grignard}, hidruro, cianuro) se adicionan directamente. Los \omterm{def:b1:electronic-effects:nucleophile}{nucleófilos}
débiles --- el agua, los alcoholes --- necesitan ayuda.

\begin{definition}[Activación electrófila]\label{def:b1:acetals-protection:activation}
La \emph{activación electrófila}\index{activacion electrofila@activación electrófila} de un grupo carbonilo
es su protonación (o su coordinación a un \omterm{def:b1:lewis-resonance-vsepr:lewis-acid}{ácido de Lewis}) en el oxígeno.
El catión \ce{C=OH+}, cuya carga positiva comparte el carbono a través de
la \omterm{def:b1:lewis-resonance-vsepr:resonance}{forma resonante} \ce{C+-OH}, es atacado incluso por \omterm{def:b1:electronic-effects:nucleophile}{nucleófilos}
débiles.
\end{definition}

\begin{omfigure}
\setchemfig{atom sep=2.1em}
\schemestart
\chemfig{C(-[:150]R)(-[:210]H)=O}
\arrow{->[\ce{H+}]}
\chemfig{C(-[:150]R)(-[:210]H)=O^{+}-H}
\arrow{<->}
\chemfig{C^{+}(-[:150]R)(-[:210]H)-O-H}
\schemestop

{\small Activación de un aldehído por un ácido: la protonación del
oxígeno produce un catión cuya \omterm{def:b1:lewis-resonance-vsepr:resonance}{forma resonante} lleva la carga en el
carbono.}
\end{omfigure}

\section{Hidratos y hemiacetales}

El agua se adiciona de forma reversible a los aldehídos y a las cetonas,
\ce{R2C=O + H2O <=>
R2C(OH)2}; el hidrato suele ser minoritario en las cetonas y puede ser
mayoritario en el metanal. Un alcohol se adiciona del mismo modo.

\begin{definition}[Hemiacetal y acetal]\label{def:b1:acetals-protection:hemiacetal}
Un \emph{hemiacetal}\index{hemiacetal} es un compuesto en el que un
carbono lleva a la vez un grupo OH y un grupo OR, formado por la adición
de un alcohol a un grupo carbonilo. Un \emph{acetal}\index{acetal} es un
compuesto en el que un carbono lleva dos grupos OR, formado a partir de
un hemiacetal y de una segunda molécula de alcohol, con pérdida de agua.
\end{definition}

Los \omterm{def:b1:acetals-protection:hemiacetal}{hemiacetales} de cadena abierta suelen ser minoritarios en su
equilibrio con el aldehído y el alcohol; cuando el OH y el C=O pertenecen
a la misma molécula y puede cerrarse un anillo de cinco o seis miembros,
predomina el \omterm{def:b1:acetals-protection:hemiacetal}{hemiacetal} cíclico. Es el caso de la glucosa, y del más
sencillo 5-hidroxipentanal.

\begin{omfigure}
\setchemfig{atom sep=2em}
\schemestart
\chemfig{HO-[:30]-[:-30]-[:30]-[:-30]-[:30]C(=[:90]O)-[:-30]H}
\arrow{<=>}
\chemfig{*6(O-(-[:-30]OH)-----)}
\schemestop

\medskip

{\small El 5-hidroxipentanal y su \omterm{def:b1:acetals-protection:hemiacetal}{hemiacetal} cíclico, un anillo de seis
miembros que contiene el oxígeno. El OH del carbono 5 se ha adicionado al
aldehído del carbono 1; predomina la forma cíclica, como en la glucosa.}
\end{omfigure}

\section{Los acetales}

\begin{proposition}[Mecanismo de formación de los acetales]\label{prop:b1:acetals-protection:acetal-mechanism}
En medio ácido, un aldehído o una cetona y un alcohol dan el \omterm{def:b1:acetals-protection:hemiacetal}{hemiacetal} y
después el \omterm{def:b1:acetals-protection:hemiacetal}{acetal}, según las etapas: protonación del C=O; adición del
alcohol; pérdida de un protón (\omterm{def:b1:acetals-protection:hemiacetal}{hemiacetal}); protonación del OH; pérdida
de agua, que da un \omterm{def:b1:electronic-effects:intermediates}{carbocatión} estabilizado por el grupo OR restante (un
ion oxocarbenio); adición de un segundo alcohol; pérdida de un protón.
Todas las etapas son reversibles.
\end{proposition}

\begin{proof}
Cada etapa es una transferencia de protón ácido--base o la adición (o la
pérdida) de un \omterm{def:b1:electronic-effects:nucleophile}{nucleófilo} a (o desde) un carbono \omterm{def:b1:electronic-effects:nucleophile}{electrófilo}, como se
describe en el \cref{ch:b1:electronic-effects}. La pérdida de agua es
posible porque el catión formado, $\mathrm{R_2C^+{-}OR}$, tiene una \omterm{def:b1:lewis-resonance-vsepr:resonance}{forma
resonante} $\mathrm{R_2C{=}O^+R}$ en la que todos los átomos tienen octeto
(\cref{prop:b1:electronic-effects:carbocation-order}). El ácido se
consume en las protonaciones y se regenera en las desprotonaciones: es un
\omterm{def:b1:elementary-steps:catalyst}{catalizador}.
\end{proof}

\begin{omfigure}
\setchemfig{atom sep=1.9em}
\schemestart
\chemfig{C(-[:150]R)(-[:210]H)=O}
\arrow{->[\ce{H+}][]}
\chemfig{C^{+}(-[:150]R)(-[:210]H)-OH}
\arrow{->[$\mathrm{R'OH}$][\ce{-H+}]}
\chemfig{C(-[:150]R)(-[:210]H)(-[2]OR')-OH}
\schemestop

\medskip
\schemestart
\chemfig{C(-[:150]R)(-[:210]H)(-[2]OR')-OH}
\arrow{->[\ce{H+}][\ce{-H2O}]}
\chemfig{C^{+}(-[:150]R)(-[:210]H)-OR'}
\arrow{->[$\mathrm{R'OH}$][\ce{-H+}]}
\chemfig{C(-[:150]R)(-[:210]H)(-[2]OR')-OR'}
\schemestop
\par\vspace{6pt}

{\small La formación de un \omterm{def:b1:acetals-protection:hemiacetal}{acetal}, en dos líneas: primero el \omterm{def:b1:acetals-protection:hemiacetal}{hemiacetal}
(activación, adición del alcohol, pérdida de un protón) y después el
\omterm{def:b1:acetals-protection:hemiacetal}{acetal} (protonación del OH, pérdida de agua hasta el catión estabilizado,
adición de un segundo alcohol, pérdida de un protón).}
\end{omfigure}

La reacción global,
$\mathrm{RCHO} + 2\,\mathrm{R'OH} \rightleftharpoons \mathrm{RCH(OR')_2} + \ce{H2O}$,
tiene una constante de equilibrio modesta, y el agua es un producto.

\begin{proposition}[Desplazar la acetalización]\label{prop:b1:acetals-protection:dean-stark}
Si el agua formada se elimina de forma continua, la acetalización
continúa hasta que se agota el compuesto carbonílico (o el alcohol).
\end{proposition}

\begin{proof}
El \omterm{def:b1:extent-q-and-k:quotient}{cociente de reacción} $Q = [\text{acetal}][\ce{H2O}]/([\text{aldehído}]
[\mathrm{R'OH}]^2)$ se mantiene por debajo de $K$ mientras se elimina el
agua: según el \cref{ch:b1:extent-q-and-k}, la reacción sigue avanzando,
sea cual sea el valor de $K$. Con un diol (el etano-1,2-diol), el \omterm{def:b1:acetals-protection:hemiacetal}{acetal}
es cíclico, y el equilibrio es además más favorable, pues una molécula de
diol sustituye a dos de alcohol.
\end{proof}

\begin{omfigure}
\begin{tikzpicture}[font=\small]
  \pic at (0,0) {heatingmantle};
  \pic at (0,0) {rbflask={0.4}{omDef!14}};
  \pic (d) at (0,3.0) {deanstark={0.35}{omDef!35}{orange!15}};
  \pic at (0,6.4) {condenser};
  \draw[->] (2.4,4.0) node[right, align=left] {colector: el agua (capa inferior)\\ya no vuelve; el tolueno\\rebosa de vuelta al matraz} -- (0.8,3.8);
  \draw[->] (-2.2,1.0) node[left, align=right] {aldehído, diol, tolueno,\\catalizador ácido, ebullición} -- (-0.7,0.9);
  \draw[->] (1.8,8.0) node[right] {refrigerante} -- (0.4,7.8);
\end{tikzpicture}

{\small Un aparato de Dean--Stark. Los vapores de tolueno y de agua se
condensan y caen en el colector graduado; el agua, más densa y no
\omterm{def:b1:intermolecular-forces-solvents:miscibility}{miscible} con el tolueno, se hunde y se queda, y el tolueno rebosa de
vuelta al matraz. El volumen de agua recogido mide el \omterm{def:b1:extent-q-and-k:extent}{avance de la
reacción}.}
\end{omfigure}

\begin{proposition}[Estabilidad de los acetales]\label{prop:b1:acetals-protection:acetal-stability}
Los \omterm{def:b1:acetals-protection:hemiacetal}{acetales} son estables en medio neutro y básico y frente a los
\omterm{def:b1:electronic-effects:nucleophile}{nucleófilos}, los \omterm{def:b1:organomagnesium:organometallic}{reactivos de Grignard} y los reactivos de tipo hidruro;
el ácido acuoso los hidroliza de nuevo al compuesto carbonílico y el
alcohol.
\end{proposition}

\begin{proof}
El carbono de un \omterm{def:b1:acetals-protection:hemiacetal}{acetal} no lleva ningún \omterm{def:b1:electronic-effects:nucleophile}{grupo saliente} que una base o un
\omterm{def:b1:electronic-effects:nucleophile}{nucleófilo} puedan expulsar: \ce{RO-}, una \omterm{def:b1:predominant-reaction:strong-acid}{base fuerte}, no se va, y ya no
queda ningún C=O que atacar. En ácido acuoso, el mecanismo anterior
transcurre en sentido inverso: el gran exceso de agua hace que $Q < K$
para la reacción inversa.
\end{proof}

\section{Proteger y desproteger}

\begin{definition}[Grupo protector]\label{def:b1:acetals-protection:protecting-group}
Un \emph{grupo protector}\index{grupo protector} es un grupo que se
introduce en una molécula para enmascarar temporalmente una función, de
modo que no reaccione en una etapa dirigida a otra parte de la molécula.
Su introducción es la \emph{protección}\index{proteccion@protección}, y su eliminación,
la \emph{desprotección}\index{desproteccion@desprotección}.
\end{definition}

\begin{method}[Proteger, reaccionar, desproteger]\label{met:b1:acetals-protection:protect}
\begin{enumerate}
  \item Se identifica la función que reaccionaría, o que destruiría el
        reactivo, en la etapa prevista.
  \item Se elige un \omterm{def:b1:acetals-protection:protecting-group}{grupo protector} estable en esa etapa y que pueda
        eliminarse en condiciones que el resto de la molécula tolere
        (para un aldehído o una cetona frente a bases, \omterm{def:b1:electronic-effects:nucleophile}{nucleófilos} o
        hidruros: un \omterm{def:b1:acetals-protection:hemiacetal}{acetal} cíclico).
  \item Se protege; se lleva a cabo la etapa prevista; se desprotege.
  \item Se cuenta el coste: dos etapas más, cada una con su \omterm{def:b1:extent-q-and-k:fractional-extent}{rendimiento}.
\end{enumerate}
\end{method}

\begin{example}[Una cetona que estorba]\label{ex:b1:acetals-protection:ketoester}
Para adicionar un \omterm{def:b1:organomagnesium:organometallic}{reactivo de Grignard} al grupo aldehído de una molécula
que lleva también una cetona, no existe ninguna elección sencilla de
\omterm{def:b1:acetals-protection:hemiacetal}{acetal} (los aldehídos forman \omterm{def:b1:acetals-protection:hemiacetal}{acetales} con más facilidad que las cetonas,
de modo que se protegería el grupo equivocado); hay que cambiar el orden
de las etapas o los reactivos. La \omterm{def:b1:acetals-protection:protecting-group}{protección} es una estrategia, no un
reflejo.
\end{example}

\section{Ejercicios}

\begin{exercise}[$\star$]\label{exo:b1:acetals-protection:1}
Identifíquense los carbonos de \omterm{def:b1:acetals-protection:hemiacetal}{hemiacetal}, \omterm{def:b1:acetals-protection:hemiacetal}{acetal}, hidrato, éter o
alcohol en:
\ce{CH3CH(OH)OCH3}, \ce{CH3CH(OCH3)2}, \ce{CH3CH(OH)2}, \ce{CH3OCH3},
\ce{(CH3)2C(OCH2CH3)2}.
\end{exercise}

\begin{exercise}[$\star$]\label{exo:b1:acetals-protection:2}
Escríbanse la ecuación global de la formación del \omterm{def:b1:acetals-protection:hemiacetal}{acetal} del propanal con
metanol y la del \omterm{def:b1:acetals-protection:hemiacetal}{acetal} cíclico de la propanona con etano-1,2-diol.
\end{exercise}

\begin{exercise}[$\star$]\label{exo:b1:acetals-protection:3}
¿Cuáles de estos reactivos dejan un \omterm{def:b1:acetals-protection:hemiacetal}{acetal} sin cambios: la disolución de
hidróxido de sodio, el bromuro de metilmagnesio, el ácido sulfúrico
diluido en agua, el borohidruro de sodio?
\end{exercise}

\begin{exercise}[$\star$]\label{exo:b1:acetals-protection:4}
Dibújese el \omterm{def:b1:acetals-protection:hemiacetal}{hemiacetal} cíclico que forma el 4-hidroxibutanal. ¿Cuál es el
tamaño del anillo?
\end{exercise}

\begin{exercise}[$\star\star$]\label{exo:b1:acetals-protection:5}
Escríbase con \omterm{def:b1:electronic-effects:curly-arrow}{flechas curvas} el mecanismo completo de la formación,
catalizada por ácido, del \omterm{def:b1:acetals-protection:hemiacetal}{acetal} dimetílico del etanal, y demuéstrese que
el ácido se regenera.
\end{exercise}

\begin{exercise}[$\star\star$]\label{exo:b1:acetals-protection:6}
Escríbase el mecanismo de la hidrólisis ácida del
2,2-dimetoxipropano. ¿Por qué se usa una gran cantidad de agua?
\end{exercise}

\begin{exercise}[$\star\star$]\label{exo:b1:acetals-protection:7}
Se sigue con un aparato de Dean--Stark la acetalización de
\qty{0.250}{mol} de una cetona. ¿Qué volumen de agua debería recogerse con
la conversión completa (densidad de aproximadamente \qty{1.0}{g/mL})? Al
cabo de una hora se han recogido \qty{3.2}{mL}: ¿cuál es la conversión?
\end{exercise}

\begin{exercise}[$\star\star$]\label{exo:b1:acetals-protection:8}
Explíquese por qué un \omterm{def:b1:acetals-protection:hemiacetal}{acetal} cíclico con etano-1,2-diol se forma de manera
más completa que el \omterm{def:b1:acetals-protection:hemiacetal}{acetal} con dos moléculas de metanol, para el mismo
compuesto carbonílico.
\end{exercise}

\begin{exercise}[$\star\star$]\label{exo:b1:acetals-protection:9}
¿Por qué hay que eliminar (o neutralizar) el \omterm{def:b1:elementary-steps:catalyst}{catalizador} ácido antes de
tratar un compuesto protegido con un \omterm{def:b1:organomagnesium:organometallic}{reactivo de Grignard}?
\end{exercise}

\begin{exercise}[$\star\star\star$]\label{exo:b1:acetals-protection:10}
Propóngase una síntesis de la 5-hidroxi-5-metilhexan-2-ona,
\ce{(CH3)2C(OH)CH2CH2COCH3}, a partir de la 4-clorobutan-2-ona y de la
propanona, usando una \omterm{def:b1:acetals-protection:protecting-group}{protección}. Dense cada etapa, sus reactivos y la
función de cada uno.
\end{exercise}

\begin{exercise}[$\star\star\star$]\label{exo:b1:acetals-protection:11}
Demuéstrese con el \omterm{def:b1:extent-q-and-k:quotient}{cociente de reacción} que un gran exceso de metanol
también desplaza la acetalización de un aldehído, y compárese este
método con la eliminación del agua.
\end{exercise}

\begin{exercise}[$\star\star\star$]\label{exo:b1:acetals-protection:12}
Explíquese por qué el \omterm{def:b1:acetals-protection:hemiacetal}{hemiacetal} cíclico de la glucosa es atacado, en
metanol ácido, para dar un \omterm{def:b1:acetals-protection:hemiacetal}{acetal} de metilo (un metil glucósido), mientras
que los demás grupos OH de la glucosa no se convierten en éteres.
\end{exercise}

\section{Problema: un reactivo de Grignard que lleva un aldehído}

\begin{problem}[{Problema de fin de semana --- por qué el 4-bromobutanal no
puede dar un reactivo de Grignard, su protección como acetal cíclico con
un aparato de Dean--Stark, la adición de Grignard a la propanona, la
desprotección y el volumen de agua recogido}]
\label{pb:b1:acetals-protection:1}
Una química necesita 5-hidroxi-5-metilhexanal,
\ce{(CH3)2C(OH)CH2CH2CH2CHO}, y piensa prepararlo a partir del
4-bromobutanal, \ce{BrCH2CH2CH2CHO}, y de la propanona. Parte de
\qty{15.09}{g} de 4-bromobutanal. Masas molares
(\unit{g/mol}): \ce{C4H7BrO} 150.9, \ce{C2H6O2} 62.0, \ce{C6H11BrO2} 194.9,
\ce{C3H6O} 58.0, \ce{C9H18O3} 174.0, \ce{C7H14O2} 130.0, \ce{H2O} 18.0;
densidad del agua, \qty{0.997}{g/mL}.
% ledger: rho:water

\medskip
\noindent\textbf{Parte I --- El problema.}
\begin{enumerate}
  \item Escríbase el \omterm{def:b1:organomagnesium:organometallic}{reactivo de Grignard} que daría el 4-bromobutanal.
  \item ¿Por qué no puede existir ese reactivo? Escríbase la reacción que
        lo destruye.
  \item ¿Qué función hay que proteger, y frente a qué?
  \item ¿Por qué conviene un \omterm{def:b1:acetals-protection:hemiacetal}{acetal}?
  \item ¿Por qué se prefiere un \omterm{def:b1:acetals-protection:hemiacetal}{acetal} cíclico (con etano-1,2-diol)?
  \item Calcúlese la cantidad de 4-bromobutanal.
  \item Calcúlese la masa de diol para un exceso del \qty{20}{\percent}.
\end{enumerate}

\medskip
\noindent\textbf{Parte II --- La \omterm{def:b1:acetals-protection:protecting-group}{protección}.}
\begin{enumerate}[resume]
  \item Escríbase la ecuación de la acetalización con etano-1,2-diol.
  \item ¿Cuál es la función del \omterm{def:b1:elementary-steps:catalyst}{catalizador} ácido (ácido
        4-metilbencenosulfónico)?
  \item Descríbase el funcionamiento del aparato de Dean--Stark.
  \item Explíquese con $Q$ y $K$ por qué eliminar el agua desplaza la
        reacción.
  \item ¿Por qué vuelve al matraz el tolueno, y no el agua?
  \item ¿Cómo sabe la química que la reacción ha terminado?
  \item Calcúlese la masa de bromuro protegido esperada con la conversión
        completa.
\end{enumerate}

\medskip
\noindent\textbf{Parte III --- La etapa de Grignard.}
\begin{enumerate}[resume]
  \item Escríbase la formación del \omterm{def:b1:organomagnesium:organometallic}{reactivo de Grignard} a partir del
        bromuro protegido.
  \item Escríbanse su adición a la propanona y el \omterm{def:b1:alcohol-activation:alkoxide}{alcóxido} formado.
  \item ¿Por qué el tratamiento de esta etapa debe ser ligeramente básico
        o neutro (por ejemplo, cloruro de amonio acuoso) y no ácido?
  \item Calcúlese la masa de alcohol protegido para un \omterm{def:b1:extent-q-and-k:fractional-extent}{rendimiento} del
        \qty{70}{\percent} en esta etapa.
\end{enumerate}

\medskip
\noindent\textbf{Parte IV --- La \omterm{def:b1:acetals-protection:protecting-group}{desprotección}.}
\begin{enumerate}[resume]
  \item Escríbase la \omterm{def:b1:acetals-protection:protecting-group}{desprotección} en ácido acuoso.
  \item ¿Por qué el alcohol terciario sobrevive al ácido acuoso suave?
  \item El producto puede cerrar un anillo de seis miembros adicionando su
        OH al aldehído. Dibújese ese \omterm{def:b1:acetals-protection:hemiacetal}{hemiacetal} cíclico.
  \item Calcúlese la masa de hidroxialdehído con un \omterm{def:b1:extent-q-and-k:fractional-extent}{rendimiento} de
        \omterm{def:b1:acetals-protection:protecting-group}{desprotección} del \qty{90}{\percent}.
  \item ¿Qué \omterm{def:b1:extent-q-and-k:fractional-extent}{rendimiento} global desde el 4-bromobutanal alcanza la
        química?
  \item Calcúlese la masa de agua formada en la etapa de \omterm{def:b1:acetals-protection:protecting-group}{protección} con la
        conversión completa.
  \item Dese el volumen de agua recogido en el aparato de Dean--Stark con
        la conversión completa.
\end{enumerate}
\end{problem}
